\section{Свойства замкнутости}

\mytodo{Замкнутость относительно гомоморфизмов: сверить по Rambow1999, доказывается ли там замкнутость относительно всех трёх операций (гомоморфизм, обратный гомоморфизм, подстановка) или только части.}

\mytodo{Доказательство незамкнутости: конструкция языков $L_1$, $L_2$  нуждается в верификации по Seki1991.}

Свойства замкнутости MCFL во многом аналогичны свойствам контекстно-свободных языков, однако имеются и существенные отличия.

\begin{theorem}[Замкнутость MCFL]
    \label{thm:mcf_closure}
    Класс MCFL замкнут относительно следующих операций:
    \begin{enumerate}
        \item \textbf{Объединение}: если $L_1, L_2 \in MCFL$, то $L_1 \cup L_2 \in MCFL$.
        \item \textbf{Конкатенация}: если $L_1, L_2 \in MCFL$, то $L_1 \cdot L_2 \in MCFL$.
        \item \textbf{Звезда Клини}: если $L_1 \in MCFL$, то $L_1^* \in MCFL$.
        \item \textbf{Пересечение с регулярными языками}: если $L_1 \in MCFL$, а $L_2$~--- регулярный язык, то $L_1 \cap L_2 \in MCFL$.
        \item \textbf{Разность с регулярными языками}: если $L_1 \in MCFL$, а $L_2$~--- регулярный язык, то $L_1 \setminus L_2 \in MCFL$.
        \item \textbf{Гомоморфизм}: если $L_1 \in MCFL$, то $h(L_1) \in MCFL$ для любого гомоморфизма $h$.
        \item \textbf{Обратный гомоморфизм}: если $L_1 \in MCFL$, то $h^{-1}(L_1) \in MCFL$ для любого гомоморфизма $h$.
        \item \textbf{Подстановка}: если $L_1 \in m$-MCFL и для каждого символа $a \in \Sigma$ задан язык $L_a \in m$-MCFL, то результат подстановки также принадлежит $m$-MCFL.
    \end{enumerate}
\end{theorem}

\begin{proof}
    Для каждого пункта приведём конструкцию.

    \begin{enumerate}
        \item Пусть $G_1 = \langle \Sigma_1, N_1, S_1, P_1 \rangle$ и $G_2 = \langle \Sigma_2, N_2, S_2, P_2 \rangle$~--- MCFG для языков $L_1$ и $L_2$ соответственно, с непересекающимися множествами нетерминалов. Грамматика для $L_1 \cup L_2$:
              \[
              G_3 = \langle \Sigma_1 \cup \Sigma_2, N_1 \cup N_2 \cup \{S_3\}, S_3, P_1 \cup P_2 \cup \{ S_3(x_1) \leftarrow S_1(x_1),\ S_3(x_1) \leftarrow S_2(x_1) \} \rangle.
              \]

        \item Конкатенация: $G_3$ получается добавлением правила $S_3(x_1 x_2) \leftarrow S_1(x_1), S_2(x_2)$.

        \item Звезда Клини: $G_2$ получается добавлением правил $S_2(x_1) \leftarrow S_1(x_1)$, $S_2(x_1 x_2) \leftarrow S_2(x_1), S_2(x_2)$.

        \item Пересечение с регулярными языками~--- ключевое свойство для MCFL-достижимости. Построим произведение MCFG $G$ и DFA $A = \langle Q, \Sigma, q_0, \delta, F \rangle$, распознающего регулярный язык. Для каждого нетерминала $N$ ранга $k$ и каждого набора состояний $p_1,\ldots,p_{k+1} \in Q$ введём аннотированный нетерминал $N_{p_1,\ldots,p_{k+1}}$, порождающий кортеж $(w_1,\ldots,w_k)$, если существуют пути $p_1 \xrightarrow{w_1} p_2$, \ldots, $p_k \xrightarrow{w_k} p_{k+1}$ в $A$. Каждое исходное правило преобразуется: все $p_i$ в левой части соединяются с $p_j$ в правой в соответствии с тем, как переменные располагаются в строках $s_1,\ldots,s_k$. Стартовый нетерминал~--- $S_{q_0, q}$ для некоторого $q \in F$.

        \item Разность с регулярным языком: $L_1 \setminus L_2 = L_1 \cap \overline{L_2}$, а регулярные языки замкнуты относительно дополнения.
        \item Гомоморфизм: заменим каждый терминал $a$ в строках $s_i$ правил MCFG на $h(a)$. Полученная грамматика порождает $h(L_1)$.

        \item Обратный гомоморфизм и подстановка: доказательство этих свойств технически более сложно; полное доказательство можно найти в работе~\sidecite{rambow1999independent}\sidenote{В этой работе показано, что $q$-MCFL($r$) для $r \ge 2$ является substitution-closed full AFL.}.
    \end{enumerate}
\end{proof}

\begin{theorem}[Незамкнутость MCFL]
    \label{thm:mcf_nonclosure}
    Класс MCFL \emph{не} замкнут относительно следующих операций:
    \begin{enumerate}
        \item \textbf{Пересечение}: существуют $L_1, L_2 \in MCFL$, такие что $L_1 \cap L_2 \notin MCFL$.
        \item \textbf{Дополнение}: существует $L \in MCFL$, такое что $\overline{L} \notin MCFL$.
        \item \textbf{Разность}: существуют $L_1, L_2 \in MCFL$, такие что $L_1 \setminus L_2 \notin MCFL$.
    \end{enumerate}
\end{theorem}

\begin{proof}
    Для доказательства незамкнутости относительно пересечения рассмотрим языки Секи с соавторами~\sidecite{seki1991multiple}:
    \[
    L_1 = \{\,a_1^n b_1^n a_2^n b_2^n \cdots a_{m+1}^n b_{m+1}^n \mid n \ge 0\,\}, \qquad
    L_2 = \{\,a_1^n a_2^n \cdots a_{m+1}^n b_1^n b_2^n \cdots b_{m+1}^n \mid n \ge 0\,\}.
    \]
    Язык $L_1$ является $(m+1)$-MCFL($1$) (теорема~\ref{thm:hierarchy_m}), язык $L_2$ является $2$-MCFL (строится парой $(a_1^n\cdots a_{m+1}^n,\; b_1^n\cdots b_{m+1}^n)$, каждый компонент которой является КС-языком).
    В работе~\sidecite{seki1991multiple} показано, что $L_1 \cap L_2$ не является MCFL: при надлежащем выборе $m$ ранг, требуемый для порождения пересечения, превосходит $m$, а конструкция допускает масштабирование на произвольно большие значения~--- таким образом, пересечение не принадлежит ни одному из классов $m$-MCFL.

    Незамкнутость относительно дополнения следует из незамкнутости относительно пересечения по законам де Моргана:
    если бы MCFL были замкнуты относительно дополнения, то $L_1 \cap L_2 = \overline{\overline{L_1} \cup \overline{L_2}}$ также был бы MCFL, что неверно.

    Незамкнутость относительно разности следует из незамкнутости относительно пересечения, так как $L_1 \cap L_2 = L_1 \setminus \overline{L_2}$, а $\overline{L_2} \in MCFL$ не гарантировано.
\end{proof}
